\noindent\textbf{Resultado.} Se tiene \(r=1/4\). En orden \(j=1,2,3\), los valores interiores, en kelvin de exceso, son
\[
(U_1^1,U_2^1,U_3^1)=(2+2\sqrt2,\ 4+2\sqrt2,\ 2+2\sqrt2),
\]
\[
(U_1^2,U_2^2,U_3^2)=\left(2+\frac{3\sqrt2}{2},\ 3+2\sqrt2,\ 2+\frac{3\sqrt2}{2}\right).
\]
En el segundo paso, los valores exactos son aproximadamente \((4.155542,\ 5.876824,\ 4.155542)\,\mathrm K\). El máximo error nodal absoluto es \(0.04839643\,\mathrm K\), en el nodo central.

\noindent\textit{Desarrollo.} El cociente de pasos es adimensional:
\[
r=\frac{(1\,\mathrm{cm^2/s})(0.25\,\mathrm s)}{(1\,\mathrm{cm})^2}=\frac14.
\]
Los datos interiores iniciales, en kelvin de exceso, son
\[
(U_1^0,U_2^0,U_3^0)=(4\sqrt2,8,4\sqrt2).
\]
El esquema toma la forma
\[
U_j^{n+1}=\frac12U_j^n+\frac14U_{j-1}^n+\frac14U_{j+1}^n.
\]
En cada paso todos los valores del segundo miembro pertenecen al mismo nivel anterior. Para el primero,
\[
U_1^1=\tfrac12(4\sqrt2)+\tfrac14(0+8)=2+2\sqrt2,
\]
\[
U_2^1=\tfrac12(8)+\tfrac14(4\sqrt2+4\sqrt2)=4+2\sqrt2,
\]
y \(U_3^1=U_1^1\) por simetría. Para el segundo, usando únicamente ese vector ya completado,
\[
U_1^2=\tfrac12(2+2\sqrt2)+\tfrac14(0+4+2\sqrt2)
=2+\frac{3\sqrt2}{2},
\]
\[
U_2^2=\tfrac12(4+2\sqrt2)+\tfrac14[2(2+2\sqrt2)]
=3+2\sqrt2,
\]
y \(U_3^2=U_1^2\). El vector aproximado final es
\[
(4.12132034,\ 5.82842712,\ 4.12132034)\,\mathrm K.
\]
En \(t=0.5\,\mathrm s\), el factor temporal exacto es \(e^{-\pi^2/32}\). Por tanto, el vector de referencia es
\[
e^{-\pi^2/32}(4\sqrt2,8,4\sqrt2)\,\mathrm K
\approx(4.15554179,\ 5.87682355,\ 4.15554179)\,\mathrm K.
\]
Los errores absolutos son aproximadamente
\[
(0.03422144,\ 0.04839643,\ 0.03422144)\,\mathrm K,
\]
por lo que el máximo solicitado es el central. El esquema atenúa este modo algo más que la solución exacta durante los dos pasos calculados.

La comparación concierne a esos nodos, ese tiempo y ese dato inicial. No establece un orden general de convergencia ni una garantía de precisión para tiempos o perfiles diferentes.
